\chapter{Dynamic Programming}\label{ch:dynamic}

\begin{center}\itshape
An optimal policy has the property that whatever the initial state and initial decision are,\\
the remaining decisions must constitute an optimal policy with regard to the state\\
resulting from the first decision.\\
\upshape(Richard Bellman, \emph{Dynamic Programming}, 1957, Chapter~III, \S\,3)
\end{center}

In the autumn of 1950 Richard Bellman was at the RAND Corporation in Santa Monica, working for the US Air Force on what he called multistage decision
processes, and he needed a name for this work. More than thirty years later he told the story in his autobiography. The Secretary of Defense, he wrote,
could not bear the word ``research'', and RAND's sponsors had to be kept happy; ``programming'' sounded like planning and scheduling, and
``dynamic'' had a precise meaning in physics and, he observed, could not be used in a pejorative sense. The combination, in his words, was
``something not even a Congressman could object to''. (Bellman's memory was not quite exact: the Secretary he named, Charles Wilson, took office only in 1953. Stuart Dreyfus, his collaborator,
reproduces the whole passage in a short article listed at the end of the chapter.) The name stuck, and it describes nothing in particular, which is perhaps why it
fits so many things.

What it describes, in this book, is a single idea: \emph{solve every subproblem once, and in the right order}. A recursive solution that calls itself on
overlapping subproblems can take exponential time, as the naive recursion for the Fibonacci numbers does (Exercise~\ref{ex:fibcalls}); if the distinct subproblems
are few, storing their answers in a table makes it polynomial. The art is in choosing the subproblems. We shall see the classical examples, each one a small gem:
the distance between two words, and how to compute an optimal alignment in linear space (Hirschberg); the knapsack; the best order in which to multiply matrices;
optimal search trees, and the beautiful theorem of Knuth and Yao that speeds them up from cubic to quadratic time; the most likely hidden explanation of a sequence of
observations (Viterbi); the travelling salesman in time $2^nn^2$ instead of $n!$; and problems on trees.

\begin{think}
How many single-letter edits (insert a letter, delete a letter, replace a letter by another) turn \texttt{kitten} into \texttt{sitting}? Three are easy to find. Can you
\emph{prove} that two are not enough?
\end{think}

\section{The principle}\label{sec:dpprinciple}

A dynamic programme has three ingredients: a family of \emph{subproblems}, indexed by a finite set $\mathcal S$; a \emph{recurrence} expressing the answer to each subproblem
through the answers to smaller ones; and an \emph{order} on $\mathcal S$ in which every subproblem comes after those it uses. The subproblems and the ``uses'' relation
form a directed acyclic graph, and the order is a topological order of it (Theorem~\ref{thm:cycle}). Filling a table in that order computes every answer exactly once,
and correctness is a proof by induction along the order, in which the main step is always the same: \emph{an optimal solution decomposes into a first decision and optimal
solutions of the subproblems this decision leaves}. This is Bellman's principle of optimality. The running time is the sum over subproblems of the work done by the
recurrence, and the space is the size of the table, which can often be reduced by forgetting answers that will not be used again.

Equivalently, one can keep the recursive program and \emph{memoise} it, storing each answer the first time it is computed; then the order is found automatically.
The Fibonacci numbers, with subproblems $F_0,\dots,F_n$ and recurrence $F_k=F_{k-1}+F_{k-2}$, need $n-1$ additions instead of the $2F_{n+1}-1$ calls of the naive recursion.

\section{Edit distance}\label{sec:edit}

Let $a=a_1\cdots a_m$ and $b=b_1\cdots b_n$ be words. The \emph{edit distance} (or Levenshtein distance) $\mathrm{ed}(a,b)$ is the least number of operations turning $a$ into $b$,
each operation inserting a letter, deleting one, or replacing one by another. An \emph{alignment} writes the two words one above the other with gaps, for instance
\begin{center}
\texttt{k i t t e n -}\\
\texttt{s i t t i n g}
\end{center}
and costs the number of columns whose two symbols differ. Each alignment of cost $c$ gives $c$ edits (a gap below is a deletion, a gap above an insertion, a mismatch a replacement), and conversely a
sequence of $c$ edits gives an alignment of cost at most $c$ (follow each letter of $a$ that survives to its final position; the survivors keep their order). So $\mathrm{ed}(a,b)$ is the least cost of an alignment.

\begin{theorem}[Levenshtein, 1965; Wagner and Fischer, 1974]\label{thm:editdistance}
Let $E(i,j)=\mathrm{ed}(a_1\cdots a_i,\,b_1\cdots b_j)$. Then $E(i,0)=i$, $E(0,j)=j$, and for $i,j\ge1$
\[
E(i,j)=\min\big(E(i-1,j)+1,\ E(i,j-1)+1,\ E(i-1,j-1)+[a_i\ne b_j]\big).
\]
So $\mathrm{ed}(a,b)=E(m,n)$ and an optimal alignment are computed in time $O(mn)$.
\end{theorem}
\begin{proof}
An empty word is reached by deletions only, and from it by insertions only. For $i,j\ge1$, look at the last column of an optimal alignment of $a_1\cdots a_i$ with $b_1\cdots b_j$: it is $a_i$ over a gap,
a gap over $b_j$, or $a_i$ over $b_j$. Removing it leaves an alignment of $a_1\cdots a_{i-1}$ with $b_1\cdots b_j$, of $a_1\cdots a_i$ with $b_1\cdots b_{j-1}$, or of $a_1\cdots a_{i-1}$ with $b_1\cdots b_{j-1}$, which costs at least the
corresponding $E$; so $E(i,j)$ is at least the minimum. Conversely, each term of the minimum is the cost of an alignment obtained by appending that column to an optimal alignment of the shorter words.
Filling the table row by row costs $O(1)$ per entry, and an optimal alignment is read backwards from $(m,n)$ by following, at each entry, a term achieving the minimum.
\end{proof}

\begin{figure}[ht]
\centering
\fig{ch05-editgrid}
\caption{The table $E(i,j)$ for \texttt{kitten} (rows) and \texttt{sitting} (columns). Its corner says $\mathrm{ed}=3$, which answers the opening question: two edits are not enough. The shaded
cells are an optimal path: diagonal steps are matches or replacements, the final horizontal step inserts \texttt{g}. The pale band is the middle row used by Hirschberg's algorithm.}
\label{fig:editgrid}
\end{figure}

An alignment is a path in the grid of Figure~\ref{fig:editgrid}, from the corner $(0,0)$ to $(m,n)$, by steps down, right, or diagonally down-right, and its cost is the sum of the step costs:
edit distance is a shortest path in a directed acyclic grid graph (Exercise~\ref{ex:dagsp}). The value $E(m,n)$ alone needs only two rows of the table at a time, so space $O(n)$. The alignment seems to need the whole
table, which for two genomes of a million letters is a terabyte. Hirschberg's trick avoids this.

\begin{theorem}[Hirschberg, 1975]\label{thm:hirschberg}
An optimal alignment of words of lengths $m$ and $n$ can be computed in time $O(mn)$ and space $O(m+n)$.
\end{theorem}
\begin{proof}
Let $h=\lfloor m/2\rfloor$. Compute in linear space the row $E(h,j)$, $0\le j\le n$, and also the row $E'(h,j)=\mathrm{ed}(a_{h+1}\cdots a_m,\,b_{j+1}\cdots b_n)$, by running the same recurrence on the reversed words.
Every path from $(0,0)$ to $(m,n)$ in the grid meets row $h$, since each step increases the row index by at most $1$, and the part before (respectively after) its first (respectively last) visit costs at least $E(h,j)$
(respectively $E'(h,j')$). So $\mathrm{ed}(a,b)=\min_j\big(E(h,j)+E'(h,j)\big)$, and for a $j^*$ achieving it, an optimal alignment of $a_1\cdots a_h$ with $b_1\cdots b_{j^*}$ followed by one of the remaining parts is optimal. Find
these two alignments recursively (stopping when $m\le1$, where the alignment is found directly), and concatenate.

The space is linear: two rows, plus the output, plus a recursion of depth $\lceil\log_2m\rceil$ storing a constant number of indices per level. For the time, exchange $a$ and $b$ if necessary so that $m\le n$,
and look at the recursion level by level. At depth $\ell$ each subproblem has at most $m/2^\ell+1$ rows; the subproblems of that depth have disjoint sets of rows, so there are at most $m$ of them, and their column ranges
are disjoint, of lengths $n_i$ with $\sum_in_i\le n$. Computing the two rows of a subproblem with $m_i$ rows and $n_i$ columns costs at most $c(m_i+1)(n_i+1)$ for a constant $c$, so depth $\ell$ costs at most
\[
c\Big(\frac m{2^\ell}+2\Big)\sum_i(n_i+1)\le c\Big(\frac m{2^\ell}+2\Big)(n+m)\le c\,\frac{2mn}{2^\ell}+4cn,
\]
using $m\le n$. Summing over the $\lceil\log_2m\rceil+1$ depths gives at most $4c\,mn+4cn(\log_2m+2)=O(mn)$.
\end{proof}

\begin{remark}
Quadratic time looks improvable, but Backurs and Indyk proved in 2015 that an algorithm in time $O(n^{2-\varepsilon})$ for edit distance would refute the \emph{strong exponential time hypothesis}, a
conjecture about satisfiability that Tome~C3 will meet. With weights on the operations (and gaps), the same recurrence is the Needleman--Wunsch algorithm of bioinformatics.
\end{remark}

\section{Knapsacks and matrix chains}\label{sec:knapsack}

A burglar can carry a total weight $W$. Object $i$, for $1\le i\le n$, has weight $w_i$ and value $v_i$, both positive integers. Which objects should be taken to maximise the value?

\begin{theorem}\label{thm:knapsack}
Let $V(i,c)$ be the largest value of a subset of the first $i$ objects of total weight at most $c$. Then $V(0,c)=0$ and
\[
V(i,c)=\begin{cases}V(i-1,c)&\text{if }w_i>c,\\ \max\big(V(i-1,c),\ v_i+V(i-1,c-w_i)\big)&\text{if }w_i\le c.\end{cases}
\]
The optimum $V(n,W)$ and an optimal subset are computed in time and space $O(nW)$.
\end{theorem}
\begin{proof}
An optimal subset for $(i,c)$ either omits object $i$, and is then a subset of the first $i-1$ objects of weight at most $c$, or contains it, and then the rest is a subset of the first $i-1$ of weight at most $c-w_i$.
So $V(i,c)$ is at most the right-hand side, and each term is the value of a feasible subset, so equality holds. There are $(n+1)(W+1)$ entries, each computed in constant time, and the subset is read backwards.
\end{proof}

The bound $O(nW)$ is not polynomial in the size of the input, which is about $n\log_2W$ bits: doubling the number of bits of $W$ can square the time. Such an algorithm is called \emph{pseudo-polynomial}.
Tome~C3 will show that the knapsack problem is NP-hard, so this is probably the best kind of algorithm one can hope for; and Chapter~12 will turn Exercise~\ref{ex:knapsackvalue} into an
approximation scheme.

Next, to multiply matrices $A_1A_2\cdots A_n$, where $A_i$ has $p_{i-1}$ rows and $p_i$ columns, by the schoolbook method, multiplying a $p\times q$ matrix by a $q\times r$ one costs $pqr$ multiplications,
and the total depends on the parenthesisation: with sizes $10\times100$, $100\times5$, $5\times50$, the order $(A_1A_2)A_3$ costs $5000+2500=7500$ and $A_1(A_2A_3)$ costs $25000+50000=75000$.

\begin{proposition}\label{prop:matrixchain}
Let $M(i,j)$ be the least cost of computing $A_i\cdots A_j$. Then $M(i,i)=0$ and, for $i<j$,
\[
M(i,j)=\min_{i\le k<j}\big(M(i,k)+M(k+1,j)+p_{i-1}p_kp_j\big),
\]
computed for all $i<j$ in order of increasing $j-i$, in time $O(n^3)$ and space $O(n^2)$.
\end{proposition}
\begin{proof}
The last multiplication of any computation of $A_i\cdots A_j$ multiplies $A_i\cdots A_k$ by $A_{k+1}\cdots A_j$ for some $k$, which are $p_{i-1}\times p_k$ and $p_k\times p_j$; the two factors were computed independently, each at cost at least
its $M$. Conversely, optimal computations of the two factors followed by this product achieve the term for $k$. There are $O(n^2)$ entries, each a minimum of at most $n$ terms.
\end{proof}

Exercise~\ref{ex:catalan} shows that trying all parenthesisations would be exponential. Hu and Shing found in 1982 an $O(n\log n)$ algorithm, by reducing the problem to a triangulation of a polygon.

\section{Optimal search trees and the quadrangle inequality}\label{sec:obst}

Keys $1,\dots,n$ are searched with frequencies $p_1,\dots,p_n\ge0$. A binary search tree on these keys finds key $k$ after $1+\mathrm{depth}(k)$ comparisons, so its cost is $\sum_kp_k(1+\mathrm{depth}(k))$. For a dictionary,
the frequent words should be near the root, but the tree must stay a search tree: which tree is best?

For $0\le i\le j\le n$ let $w(i,j)=p_{i+1}+\dots+p_j$ and let $c(i,j)$ be the least cost of a search tree on the keys $i+1,\dots,j$ (so $c(i,i)=0$). A tree on these keys with root $k$ has a left subtree on $i+1,\dots,k-1$ and a right subtree
on $k+1,\dots,j$, and all depths in them increase by $1$; the root itself contributes $p_k$. Hence, exactly as in Proposition~\ref{prop:matrixchain}, for $i<j$,
\begin{equation}\label{eq:obst}
c(i,j)=w(i,j)+\min_{i<k\le j}\big(c(i,k-1)+c(k,j)\big),
\end{equation}
and an optimal tree is found in time $O(n^3)$. Figure~\ref{fig:obst} shows one. In 1971 Knuth noticed that the optimal roots move monotonically, and that this saves a factor $n$. In 1980 Frances Yao explained why, with
an inequality that has since been used for dozens of other problems.

\begin{figure}[ht]
\centering
\fig{ch05-obst}
\caption{An optimal search tree for keys $A<B<\dots<G$ with frequencies $4,1,6,2,9,3,5$ (in red). Its cost is $62$. Frequent keys are near the root, but not in order of frequency: $D$, searched twice, is as high as
$A$, searched four times, because the tree must respect the order of the keys.}
\label{fig:obst}
\end{figure}

Say that a function $f$ of pairs $i\le j$ satisfies the \emph{quadrangle inequality} if
\[
f(i,j)+f(i',j')\le f(i',j)+f(i,j')\qquad\text{whenever } i\le i'\le j\le j',
\]
and that it is \emph{monotone} if $f(i',j)\le f(i,j')$ whenever $[i',j]\subseteq[i,j']$. For \eqref{eq:obst}, $w$ satisfies the quadrangle inequality with equality (both sides add the same $p$'s) and is monotone since $p\ge0$.
In what follows $c$ is defined by \eqref{eq:obst} from any such $w$; write $c_k(i,j)=w(i,j)+c(i,k-1)+c(k,j)$, and let $K(i,j)$ be the \emph{largest} $k$ achieving the minimum in \eqref{eq:obst}.

\begin{lemma}[Yao, 1980]\label{lem:yaoqi}
If $w$ satisfies the quadrangle inequality and is monotone, then so does $c$ satisfy the quadrangle inequality.
\end{lemma}
\begin{proof}
Induction on $\ell=j'-i$. If $i=i'$ or $j=j'$ there is equality, so let $i<i'\le j<j'$; the induction hypothesis covers all quadruples spanning fewer than $\ell$.

\emph{Case $i'=j$.} We must show $c(i,j)+c(j,j')\le c(i,j')$. Let $z=K(i,j')$. If $z\le j$, then $c(i,j)\le c_z(i,j)$, so
\[\begin{aligned}
c(i,j)+c(j,j')&\le w(i,j)+c(i,z-1)+\big(c(z,j)+c(j,j')\big)\\
&\le w(i,j')+c(i,z-1)+c(z,j')=c(i,j'),
\end{aligned}\]
using monotonicity of $w$ and the induction hypothesis on $z\le j\le j\le j'$ (which spans $j'-z<\ell$). If $z>j$, then $c(j,j')\le c_z(j,j')$, and symmetrically
\[\begin{aligned}
c(i,j)+c(j,j')&\le w(j,j')+\big(c(i,j)+c(j,z-1)\big)+c(z,j')\\
&\le w(i,j')+c(i,z-1)+c(z,j')=c(i,j').
\end{aligned}\]

\emph{Case $i'<j$.} Let $y=K(i',j)$ and $z=K(i,j')$. If $z\le y$, then $i<z\le j$ and $i'<y\le j'$, so
\[\begin{aligned}
c(i,j)+c(i',j')&\le c_z(i,j)+c_y(i',j')\\
&=w(i,j)+w(i',j')+c(i,z-1)+c(i',y-1)+c(z,j)+c(y,j').
\end{aligned}\]
By the hypothesis on $w$, $w(i,j)+w(i',j')\le w(i',j)+w(i,j')$, and by induction on $z\le y\le j\le j'$, $c(z,j)+c(y,j')\le c(y,j)+c(z,j')$. Regrouping, the right-hand side is at most
$c_y(i',j)+c_z(i,j')=c(i',j)+c(i,j')$. If $z>y$, use instead $c(i,j)\le c_y(i,j)$ (valid as $i<y\le j$) and $c(i',j')\le c_z(i',j')$ (valid as $i'<z\le j'$), and the induction hypothesis on
$i\le i'\le y-1\le z-1$: $c(i,y-1)+c(i',z-1)\le c(i',y-1)+c(i,z-1)$. Regrouping gives again $c_y(i',j)+c_z(i,j')$.
\end{proof}

\begin{theorem}[Knuth, 1971; Yao, 1980]\label{thm:knuthyao}
Under the hypotheses of Lemma~\ref{lem:yaoqi}, $K(i,j-1)\le K(i,j)\le K(i+1,j)$ whenever $j-i\ge2$. Consequently, computing $c(i,j)$ by trying only $k$ between $K(i,j-1)$ and $K(i+1,j)$, in order of increasing $j-i$, takes
time $O(n^2)$. In particular, optimal binary search trees are found in quadratic time.
\end{theorem}
\begin{proof}
\emph{Left inequality.} Let $k\le k'\le j-1$. The quadrangle inequality for $c$ on $k\le k'\le j-1\le j$, with $w(i,j-1)+w(i,j)+c(i,k-1)+c(i,k'-1)$ added to both sides, reads
\[
c_k(i,j-1)+c_{k'}(i,j)\le c_{k'}(i,j-1)+c_k(i,j).
\]
So if $k'$ is at least as good as $k$ for $(i,j-1)$, that is $c_{k'}(i,j-1)\le c_k(i,j-1)$, then $c_{k'}(i,j)\le c_k(i,j)$. Taking $k'=K(i,j-1)$: no $k<k'$ beats $k'$ for $(i,j)$, so the largest minimiser $K(i,j)$ is at least $k'$.

\emph{Right inequality.} Let $k'=K(i+1,j)$ and $k'<k\le j$. The quadrangle inequality on $i\le i+1\le k'-1\le k-1$, with $w(i,j)+w(i+1,j)+c(k',j)+c(k,j)$ added, reads
\[
c_{k'}(i,j)+c_k(i+1,j)\le c_{k'}(i+1,j)+c_k(i,j).
\]
Since $k$ is larger than the largest minimiser for $(i+1,j)$, $c_k(i+1,j)>c_{k'}(i+1,j)$, hence $c_k(i,j)>c_{k'}(i,j)$: $k$ is not a minimiser for $(i,j)$, and $K(i,j)\le k'$.

\emph{Time.} For $j-i=d\ge2$, computing $c(i,j)$ and $K(i,j)$ costs $O\big(K(i+1,j)-K(i,j-1)+1\big)$. Summed over $i$, with $j=i+d$, the differences telescope to $K(n-d+1,n)-K(0,d-1)\le n$, so each diagonal costs $O(n)$
and the $n$ diagonals cost $O(n^2)$.
\end{proof}

\begin{beauty}
The cubic algorithm tries every root for every interval. Knuth and Yao show that this is wasteful for a reason that has nothing to do with trees: the cost satisfies an inequality between four numbers,
$c(i,j)+c(i',j')\le c(i',j)+c(i,j')$, which says that two overlapping intervals cost no more than the union and the intersection; from it, the best choices march monotonically across the table, and a
sum that looked quadratic per diagonal telescopes to a linear one.
\end{beauty}

\section{Hidden paths: Viterbi}\label{sec:viterbi}

In 1967 Andrew Viterbi, then at the University of California at Los Angeles, published a method for decoding convolutional codes. It turned out to be a dynamic programme, and today it decodes signals in
mobile phones and satellite links, aligns speech with words, and finds genes in DNA. Its setting is a \emph{hidden Markov model}: a finite set $S$ of hidden states, an initial distribution $\pi$, transition
probabilities $a(s,s')$, and emission probabilities $b(s,o)$ of observing the symbol $o$ in state $s$. A hidden path $x_1,\dots,x_T$ and observations $o_1,\dots,o_T$ have joint probability
\[
P(x,o)=\pi(x_1)b(x_1,o_1)\prod_{t=2}^Ta(x_{t-1},x_t)\,b(x_t,o_t).
\]
Given the observations, which hidden path is the most likely? There are $|S|^T$ candidates.

\begin{theorem}[Viterbi, 1967]\label{thm:viterbi}
Let $\delta_t(s)$ be the largest value of $P(x_1\cdots x_t,o_1\cdots o_t)$ over hidden paths of length $t$ ending in $x_t=s$. Then $\delta_1(s)=\pi(s)b(s,o_1)$ and
\[
\delta_t(s')=\max_{s\in S}\delta_{t-1}(s)\,a(s,s')\,b(s',o_t).
\]
So a most likely hidden path is found in time $O(T|S|^2)$.
\end{theorem}
\begin{proof}
The probability of a path of length $t$ ending in $s'$ is the probability of its prefix of length $t-1$, ending in some $s$, times $a(s,s')b(s',o_t)$, a factor that depends only on $s$ and $s'$. Maximising over the prefix for fixed $s$
gives $\delta_{t-1}(s)a(s,s')b(s',o_t)$, then maximise over $s$. Keep for each $(t,s')$ a maximising $s$; the best path ends at a state maximising $\delta_T$ and is read backwards.
\end{proof}

Taking $-\log$ of the probabilities, this is a shortest path in the layered acyclic graph with vertices $(t,s)$, and in practice the logarithms also prevent the products from underflowing. Replacing $\max$ by a sum gives
the \emph{forward algorithm}, which computes $P(o)$ itself (Exercise~\ref{ex:forward}).

\section{The travelling salesman}\label{sec:heldkarp}

Cities $1,\dots,n$ are at distances $d(i,j)\ge0$, not necessarily symmetric. A tour visits every city once and returns. There are $(n-1)!$ tours from city $1$, and $30!$ is about $2.7\times10^{32}$.

\begin{theorem}[Bellman, 1962; Held and Karp, 1962]\label{thm:heldkarp}
For $S\subseteq\{2,\dots,n\}$ and $j\in S$, let $C(S,j)$ be the least length of a path that starts at $1$, visits exactly the cities of $S\cup\{1\}$, and ends at $j$. Then $C(\{j\},j)=d(1,j)$ and, for $|S|\ge2$,
\[
C(S,j)=\min_{i\in S\setminus\{j\}}\big(C(S\setminus\{j\},i)+d(i,j)\big).
\]
The shortest tour has length $\min_jC(\{2,\dots,n\},j)+d(j,1)$, computed in time $O(2^nn^2)$ and space $O(2^nn)$.
\end{theorem}
\begin{proof}
In a path as in the definition, with $|S|\ge2$, the city $i$ visited just before $j$ is in $S\setminus\{j\}$, and the part up to $i$ is a path from $1$ visiting exactly $(S\setminus\{j\})\cup\{1\}$ and ending at $i$; it is at least
$C(S\setminus\{j\},i)$ long. Conversely, extending an optimal such path by the edge $ij$ achieves the term for $i$. A tour is such a path for $S=\{2,\dots,n\}$ closed by an edge back to $1$. Processing the sets by increasing size, there
are at most $2^{n-1}(n-1)$ pairs $(S,j)$, each a minimum over fewer than $n$ terms. Encoding $S$ as a bit mask makes $S\setminus\{j\}$ a single operation.
\end{proof}

For $n=30$ this is about $10^{12}$ steps instead of $10^{32}$: feasible in a day instead of never. Whether the travelling salesman problem can be solved in time $O(1.999^n)$ is a famous open question; for deciding whether a
Hamiltonian cycle exists, Bj\"orklund broke the $2^n$ barrier for undirected graphs in 2010.

\section{Dynamic programming on trees}\label{sec:treedp}

Trees are the natural home of dynamic programming: the subtrees of a rooted tree are subproblems, and they overlap in a perfectly controlled way. A typical example: in a tree with vertex weights $x(v)\ge0$, find
an \emph{independent set} (no two of its vertices adjacent) of largest weight. In a general graph this is NP-hard.

\begin{proposition}\label{prop:treemis}
Root the tree. For a vertex $v$ with children $c_1,\dots,c_r$, let $\mathrm{in}(v)$ and $\mathrm{out}(v)$ be the largest weights of independent sets of the subtree of $v$ that contain $v$, respectively do not contain $v$. Then
\[
\mathrm{in}(v)=x(v)+\sum_i\mathrm{out}(c_i),\qquad \mathrm{out}(v)=\sum_i\max\big(\mathrm{in}(c_i),\mathrm{out}(c_i)\big),
\]
and the maximum weight is $\max(\mathrm{in},\mathrm{out})$ at the root, computed in linear time by a postorder traversal.
\end{proposition}
\begin{proof}
An independent set $I$ of the subtree of $v$ is the union of its traces on the subtrees of the children, plus possibly $v$; there is no edge between different subtrees, so these traces can be chosen independently, except that
if $v\in I$, no child is in $I$. This gives both formulas, by maximising each trace separately. Each vertex is handled once with work proportional to its number of children.
\end{proof}

\begin{remark}
The same method works on graphs that are ``tree-like'': if a graph has a tree decomposition of width $k$, independent set, colouring and many other NP-hard problems are solved in time $f(k)\cdot n$. Courcelle's theorem (1990)
says that every property expressible in monadic second-order logic can be decided this way, a striking meeting of logic and algorithms.
\end{remark}

\section{Exercises}

\begin{exercise}[$\star$]\label{ex:lcs}
A \emph{common subsequence} of $a$ and $b$ is a word obtained from both by deleting letters. Give a recurrence for the length $L(i,j)$ of a longest common subsequence of $a_1\cdots a_i$ and $b_1\cdots b_j$. Show that the least number of
insertions and deletions (without replacements) turning $a$ into $b$ is $m+n-2L(m,n)$.
\end{exercise}

\begin{exercise}[$\star$]\label{ex:coins}
With coins of values $1$, $3$ and $4$, show that the greedy method (take the largest coin that fits) does not always use the fewest coins. Give an $O(nV)$ algorithm for the fewest coins paying $V$ with $n$ coin values, and an algorithm
counting the ways to pay $V$.
\end{exercise}

\begin{exercise}[$\star\star$]\label{ex:lis}
Find a longest increasing subsequence of a sequence of $n$ numbers in time $O(n\log n)$. (Keep, for each length $\ell$, the smallest possible last element of an increasing subsequence of length $\ell$.)
\end{exercise}

\begin{exercise}[$\star$]\label{ex:catalan}
Show that the number of parenthesisations of a product of $n$ factors is the Catalan number $C_{n-1}=\frac1n\binom{2n-2}{n-1}$, and that it is at least $2^{n-2}$ for $n\ge2$.
\end{exercise}

\begin{exercise}[$\star$]\label{ex:knapsackvalue}
Solve the knapsack problem in time $O(n^2v_{\max})$, where $v_{\max}=\max_iv_i$, by computing for each value $v$ the least weight of a subset of value exactly $v$.
\end{exercise}

\begin{exercise}[$\star\star$]\label{ex:linebreak}
\emph{Breaking paragraphs into lines.} Words of lengths $\ell_1,\dots,\ell_n$ must be set in lines of width $L$, in order, with one space between words on a line. A line holding words $i$ to $j$ has slack $s=L-\sum_{k=i}^j\ell_k-(j-i)$, which must
be nonnegative, and costs $s^2$, except the last line, which is free. Minimise the total cost in time $O(n^2)$. Show that the greedy method, filling each line as much as possible, is not optimal.
\end{exercise}

\begin{exercise}[$\star\star$]\label{ex:bitonic}
\emph{Bitonic tours.} Points $p_1,\dots,p_n$ in the plane have increasing abscissas. A tour is \emph{bitonic} if it goes from $p_1$ to $p_n$ strictly rightwards and comes back strictly leftwards. Find the shortest bitonic tour in time $O(n^2)$.
\end{exercise}

\begin{exercise}[$\star$]\label{ex:treevc}
Find a minimum vertex cover of a tree in linear time. Show that in a tree its size equals the size of a maximum matching.
\end{exercise}

\begin{exercise}[$\star$]\label{ex:forward}
\emph{The forward algorithm.} With the notation of Theorem~\ref{thm:viterbi}, let $\alpha_t(s)=P(o_1\cdots o_t,\,x_t=s)$. Show that $\alpha_t(s')=\sum_s\alpha_{t-1}(s)a(s,s')b(s',o_t)$ and deduce $P(o_1\cdots o_T)$ in time $O(T|S|^2)$.
\end{exercise}

\begin{exercise}[$\star\star$]\label{ex:hirschspace}
Show that the length of a longest common subsequence can be computed in space $O(\min(m,n))$, and that Hirschberg's method also gives the subsequence itself in time $O(mn)$ and linear space.
\end{exercise}

\section{Hints}

\begin{hint}{ex:lcs}
Compare $a_i$ and $b_j$. For the second part, the letters that are neither deleted nor inserted form a common subsequence.
\end{hint}
\begin{hint}{ex:coins}
Pay $6$. For counting, process the coin values one at a time, to count each multiset once.
\end{hint}
\begin{hint}{ex:lis}
The array of smallest last elements is increasing; a new number replaces the first entry that is at least as large.
\end{hint}
\begin{hint}{ex:catalan}
Split at the last multiplication; the Catalan numbers satisfy $C_k=\sum_{i=0}^{k-1}C_iC_{k-1-i}$.
\end{hint}
\begin{hint}{ex:knapsackvalue}
The total value is at most $nv_{\max}$.
\end{hint}
\begin{hint}{ex:linebreak}
Let $f(j)$ be the least cost of setting words $1,\dots,j$, with word $j$ ending a line. For greed, try $L=10$ and words of lengths $5,4,4,10$.
\end{hint}
\begin{hint}{ex:bitonic}
Consider two disjoint rightward paths from $p_1$ covering $p_1,\dots,p_j$ and ending at $p_i$ and $p_j$, $i<j$. Where is $p_{j-1}$?
\end{hint}
\begin{hint}{ex:treevc}
Two values per vertex, as in Proposition~\ref{prop:treemis}. For the matching, match leaves to their parents greedily.
\end{hint}
\begin{hint}{ex:forward}
Sum instead of maximising, by the law of total probability.
\end{hint}
\begin{hint}{ex:hirschspace}
Keep two rows; for the subsequence, split $a$ in the middle and combine forward and backward rows.
\end{hint}

\section{Solutions}

\begin{solution}{ex:lcs}
\idea Look at the last letters: either they match and can be used together, or one of them is useless.
\writeup $L(i,0)=L(0,j)=0$; if $a_i=b_j$, then $L(i,j)=L(i-1,j-1)+1$, otherwise $L(i,j)=\max(L(i-1,j),L(i,j-1))$. Proof: a common subsequence of $a_1\cdots a_i$ and $b_1\cdots b_j$ does not use $a_i$, or does not use $b_j$, or uses both as its
last letter, which requires $a_i=b_j$; when $a_i=b_j$, a longest one can be assumed to end with this letter (if it ends with $a_i$ matched to an earlier copy, rematch it to $b_j$; if it uses neither, append the letter). For the second part, the letters of $a$ kept by an
insertion--deletion sequence form a common subsequence of length $k$, and the sequence makes $m-k$ deletions and $n-k$ insertions; conversely, a common subsequence of length $k$ gives such a sequence. So the minimum is $m+n-2L(m,n)$.
\end{solution}

\begin{solution}{ex:coins}
\idea Greed commits to the big coin; dynamic programming considers every last coin.
\writeup For $6$, greed takes $4+1+1$, three coins, while $3+3$ uses two. Let $N(v)$ be the fewest coins paying $v$: $N(0)=0$ and $N(v)=1+\min\{N(v-c): c\le v\}$ over the coin values, since the last coin is some $c$ and the rest pays $v-c$ optimally. This is $O(nV)$. To count
multisets, let $P_i(v)$ be the number of ways using only the first $i$ values: $P_i(v)=P_{i-1}(v)+P_i(v-c_i)$ (either value $c_i$ is not used at all, or remove one coin $c_i$), with $P_i(0)=1$; again $O(nV)$.
\end{solution}

\begin{solution}{ex:lis}
\idea Among increasing subsequences of a given length, the one with the smallest last element is the most useful for the future.
\writeup Scan the sequence, keeping $T[\ell]$, the smallest last element of an increasing subsequence of length $\ell$ of the prefix seen so far. $T$ is increasing: from a subsequence of length $\ell+1$ ending in $T[\ell+1]$, dropping its last element leaves one of length
$\ell$ ending in something smaller. A new number $x$ extends exactly the subsequences ending below $x$; so if $\ell$ is the largest index with $T[\ell]<x$, the new subsequence of length $\ell+1$ ending in $x$ improves $T[\ell+1]$ (which was $\ge x$, or undefined), and no other entry
changes. Finding $\ell$ is a binary search, $O(\log n)$. The answer is the final length of $T$; storing, for each element, a pointer to the element that was $T[\ell]$ when it was inserted recovers a longest subsequence.
\end{solution}

\begin{solution}{ex:catalan}
\idea A parenthesisation is a binary tree with $n$ leaves, and the root splits it in two.
\writeup Let $P_n$ be the number for $n$ factors. The last multiplication splits the factors into the first $k$ and the last $n-k$, $1\le k\le n-1$, parenthesised independently: $P_1=1$, $P_n=\sum_{k=1}^{n-1}P_kP_{n-k}$. Setting $C_{k}=P_{k+1}$ gives
the Catalan recurrence with $C_0=1$, whose solution is $\frac1{k+1}\binom{2k}k$ (Tome~C1, Corollary~3.12, or the generating function $C(x)=1+xC(x)^2$). For the bound, induction: $P_2=1$, and $P_n\ge P_1P_{n-1}+P_{n-1}P_1=2P_{n-1}$ for $n\ge3$.
\end{solution}

\begin{solution}{ex:knapsackvalue}
\idea Exchange the roles of weight and value: tabulate by value, which is small when values are small integers.
\writeup Let $G(i,v)$ be the least weight of a subset of the first $i$ objects with value exactly $v$ ($\infty$ if none), for $0\le v\le nv_{\max}$. Then $G(0,0)=0$, $G(0,v)=\infty$ for $v>0$, and $G(i,v)=\min\big(G(i-1,v),\,w_i+G(i-1,v-v_i)\big)$, the second term only if $v_i\le v$,
by the argument of Theorem~\ref{thm:knapsack}. The answer is the largest $v$ with $G(n,v)\le W$. The table has $(n+1)(nv_{\max}+1)$ entries.
\end{solution}

\begin{solution}{ex:linebreak}
\idea Decide where the last line starts; everything before it is an optimal smaller paragraph.
\writeup Let $f(0)=0$ and, for $j\ge1$, $f(j)=\min_i\big(f(i-1)+\mathrm{cost}(i,j)\big)$ over the $i\le j$ such that words $i..j$ fit on a line, where $\mathrm{cost}(i,j)=s^2$, or $0$ if $j=n$. The proof is the usual one: in an optimal setting of words $1..j$ the last line holds words $i..j$ for some $i$, and the
lines before are a setting of $1..i-1$. With prefix sums, each cost is $O(1)$: time $O(n^2)$. For greed, take $L=10$ and lengths $5,4,4,10$: greed sets ``$5\ 4$'' (slack $0$), then ``$4$'' alone (slack $6$, cost $36$), then the last word, total $36$;
the setting ``$5$'' (slack $5$), ``$4\ 4$'' (slack $1$), last word, costs $25+1=26$. Greed is blind to the line it is preparing for the next word. Knuth and Plass's algorithm in \TeX{} is this dynamic programme with a richer cost; it set the line breaks of this book.
\end{solution}

\begin{solution}{ex:bitonic}
\idea Sweep from left to right keeping the two ends of the two strands of the tour; only the leftmost end matters, the other one is the current point.
\writeup For $i<j$ let $P(i,j)$ be the least total length of two vertex-disjoint paths, each going strictly rightwards, both starting at $p_1$ (which they share), together covering $p_1,\dots,p_j$, one ending at $p_i$ and the other at $p_j$. The
path ending at $p_i$ uses only points of index at most $i$. If $i<j-1$, the point $p_{j-1}$ must therefore be on the other path, just before $p_j$: $P(i,j)=P(i,j-1)+|p_{j-1}p_j|$. If $i=j-1$, the point before $p_j$ is some $p_k$, $k<j-1$, and removing $p_j$ leaves a pair ending at $p_k$ and $p_{j-1}$:
$P(j-1,j)=\min_{k<j-1}\big(P(k,j-1)+|p_kp_j|\big)$, with $P(1,2)=|p_1p_2|$. The shortest bitonic tour has length $P(n-1,n)+|p_{n-1}p_n|$, since $p_{n-1}$ and $p_n$ are consecutive on it. There are $O(n^2)$ entries; the first kind costs $O(1)$ and the $n$ entries of the second kind $O(n)$ each.
\end{solution}

\begin{solution}{ex:treevc}
\idea The edge from a leaf to its parent must be covered, and the parent covers more.
\writeup With $\mathrm{in}(v)=1+\sum_i\min(\mathrm{in}(c_i),\mathrm{out}(c_i))$ and $\mathrm{out}(v)=\sum_i\mathrm{in}(c_i)$ (if $v$ is not in the cover, every edge to a child forces the child in), the minimum is $\min(\mathrm{in},\mathrm{out})$ at the root, by the argument of Proposition~\ref{prop:treemis}. For the equality, a cover has at least as many vertices
as any matching has edges, since the matched edges are disjoint and each needs its own vertex. Conversely, repeatedly take a leaf $u$ with parent $v$ (in the current forest; isolated vertices are discarded), put $v$ in the cover and $uv$ in the matching, and delete $u$ and $v$. Every deleted edge is covered by $v$, and the
matching and the cover grow together. (This is K\H{o}nig's theorem for trees; it holds for all bipartite graphs, Tome~C1, Theorem~6.4.)
\end{solution}

\begin{solution}{ex:forward}
\idea Probabilities of disjoint events add; the Markov structure makes the last step a factor depending only on the last two states.
\writeup The events $\{x_{t-1}=s\}$ partition the sample space, and $P(o_1\cdots o_t,x_{t-1}=s,x_t=s')=P(o_1\cdots o_{t-1},x_{t-1}=s)\,a(s,s')\,b(s',o_t)$, by the product formula for $P(x,o)$ summed over $x_1,\dots,x_{t-2}$. Summing over $s$ gives the recurrence, with $\alpha_1(s)=\pi(s)b(s,o_1)$,
and $P(o_1\cdots o_T)=\sum_s\alpha_T(s)$.
\end{solution}

\begin{solution}{ex:hirschspace}
\idea Row $i$ of the table depends only on row $i-1$; and the middle row splits the problem in two, as in Theorem~\ref{thm:hirschberg}.
\writeup Exchanging $a$ and $b$ if needed, let $n\le m$. Compute $L(i,\cdot)$ for $i=1,\dots,m$ keeping only the previous row, of length $n+1$. For the subsequence, with $h=\lfloor m/2\rfloor$, compute the row $L(h,j)$ and the row $L'(h,j)$, the length of a longest common subsequence of
$a_{h+1}\cdots a_m$ and $b_{j+1}\cdots b_n$; a longest common subsequence of $a$ and $b$ splits into one of $a_1\cdots a_h$ with $b_1\cdots b_j$ and one of the rest, for some $j$, so $L(m,n)=\max_j\big(L(h,j)+L'(h,j)\big)$. Recurse on both halves with a maximising $j$. The time analysis is that of
Theorem~\ref{thm:hirschberg}. This is the setting of Hirschberg's original paper.
\end{solution}

\section*{Sources and further reading}

\begin{itemize}
\item \emph{The story.} R.~Bellman, \emph{Eye of the Hurricane: An Autobiography}, World Scientific, 1984, p.~159, quoted and discussed in S.~Dreyfus, \emph{Richard Bellman on the birth of dynamic programming}, Operations Research 50 (2002), 48--51.
The epigraph is from R.~Bellman, \emph{Dynamic Programming}, Princeton University Press, 1957, Chapter~III, \S\,3.
\item V.~I.~Levenshtein, \emph{Binary codes capable of correcting deletions, insertions, and reversals}, Doklady Akademii Nauk SSSR 163 (1965), 845--848 (English translation in Soviet Physics Doklady 10 (1966), 707--710); R.~A.~Wagner and M.~J.~Fischer,
\emph{The string-to-string correction problem}, Journal of the ACM 21 (1974), 168--173; D.~S.~Hirschberg, \emph{A linear space algorithm for computing maximal common subsequences}, Communications of the ACM 18 (1975), 341--343; S.~B.~Needleman and C.~D.~Wunsch,
\emph{A general method applicable to the search for similarities in the amino acid sequence of two proteins}, Journal of Molecular Biology 48 (1970), 443--453; A.~Backurs and P.~Indyk, \emph{Edit distance cannot be computed in strongly subquadratic time (unless SETH is false)},
SIAM Journal on Computing 47 (2018), 1087--1097 (conference version STOC 2015).
\item D.~E.~Knuth, \emph{Optimum binary search trees}, Acta Informatica 1 (1971), 14--25; F.~F.~Yao, \emph{Efficient dynamic programming using quadrangle inequalities}, Proceedings of the 12th ACM STOC (1980), 429--435; T.~C.~Hu and M.~T.~Shing, \emph{Computation of matrix chain products},
Part~I, SIAM Journal on Computing 11 (1982), 362--373, and Part~II, 13 (1984), 228--251.
\item A.~J.~Viterbi, \emph{Error bounds for convolutional codes and an asymptotically optimum decoding algorithm}, IEEE Transactions on Information Theory 13 (1967), 260--269; L.~R.~Rabiner, \emph{A tutorial on hidden Markov models and selected applications in speech recognition},
Proceedings of the IEEE 77 (1989), 257--286.
\item R.~Bellman, \emph{Dynamic programming treatment of the travelling salesman problem}, Journal of the ACM 9 (1962), 61--63; M.~Held and R.~M.~Karp, \emph{A dynamic programming approach to sequencing problems}, Journal of SIAM 10 (1962), 196--210; A.~Bj\"orklund, \emph{Determinant sums for
undirected Hamiltonicity}, SIAM Journal on Computing 43 (2014), 280--299 (conference version FOCS 2010).
\item B.~Courcelle, \emph{The monadic second-order logic of graphs I: recognizable sets of finite graphs}, Information and Computation 85 (1990), 12--75; D.~E.~Knuth and M.~F.~Plass, \emph{Breaking paragraphs into lines}, Software: Practice and Experience 11 (1981), 1119--1184.
\end{itemize}
